\documentclass[8pt]{beamer}

\setbeamertemplate{bibliography item}{\insertbiblabel}

\usepackage{etex}
\usecolortheme[RGB={205,173,0}]{structure} 
\usefonttheme{serif} 

%Pausing toggle
\def\mypause{
\pause
} %Remove "%" or comment out

%Style
\usepackage{indentfirst}
%\usepackage{cite} 					%Hyphened citations
%\usepackage[symbol]{footmisc}

%Viewing
\usepackage{hyperref}
%\hypersetup{
%	unicode,	% use with \texorpdfstring
%	colorlinks,
%%	citecolor=cyan,% filecolor=black,% 
%%	linkcolor=cyan, % urlcolor=black,% pdftex
%%	bookmarks=true,
%	bookmarksopen=true,
%	bookmarksopenlevel=\maxdimen,
%%	pdfpagemode=FullScreen
%%	pdfmenubar=false
%} 

%Math
\usepackage{amsmath, amssymb} 	%Needed for align environment
\usepackage{mathrsfs} 						%For \mathscr
\usepackage{bm}								%For bold greek letter and such; command is \bm
%\usepackage[mathcal]{euscript}		%For different script letters
\def\mb#1{{\bm #1}}
\def\multi#1#2{ {#1}_{\mb#2}}

%Drawing
\usepackage[all]{xy}
\usepackage{tikz}
	\usetikzlibrary{decorations.markings}
	\usetikzlibrary{arrows}

%To get equations numbered by section
\DeclareGraphicsRule{.tif}{png}{.png}{`convert #1 `dirname #1`/`basename #1 .tif`.png}
%\numberwithin{equation}{subsection}
\numberwithin{equation}{section}


%%%%%%%%%%%%%%%%%%%%%%%%%%%%%%%%%%%%%%%%%%%%%%%%%%
%%%%%%%%%%%%%%%%  Theorem environments  %%%%%%%%%%%%%%%%%%%%
% Theorem environments
\usepackage{amsthm} %Enable theorem environments
%% No-number theorems
%\theoremstyle{definition}
%\newtheorem*{definition}{Definition}
%\theoremstyle{plain}
%\newtheorem*{prop}{Proposition}
%\newtheorem*{lemma}{Lemma}
%\newtheorem*{thm}{Theorem}
%\newtheorem*{cor}{Corollary}
%\newtheorem*{example}{Example}
%% Numbered theorems
%\theoremstyle{definition}
%\newtheorem{definition}{Definition}
%\theoremstyle{plain}
%\newtheorem{proposition}{Proposition}
%\newtheorem{lemma}{Lemma}
%\newtheorem{theorem}{Theorem}
%\newtheorem{corollary}{Corollary}
%\newtheorem{conjecture}{Conjecture}
%\newtheorem{example}{Example}
%\newtheorem{result}{Result}
\newtheorem*{result*}{Result}

\usepackage{xparse}

% accent over:
\def\on#1#2{{\buildrel{\mkern2.5mu#1\mkern-2.5mu}\over{#2}}}
\def\dt#1{\on{\hbox{\bf .}}{#1}}                % (big) dot over: see�   below  
\def\under#1#2{\mathop{\null#2}\limits_{#1}}	% accent under
\mathcode`\*="702A                  % * now always complex conjugate
\def\f#1#2{{\textstyle{#1\over#2}}}	   % fraction
\def\half{{\textstyle{1\over{\raise.1ex\hbox{$\scriptstyle{2}$}}}}}

%%%%%%%%%%%%%%%%%%%%%%%%%%%%%%%%%%%%%%%%%%%%%%%%%%
%%%%%%%%%%%%%%%%  Typesetting conventions  %%%%%%%%%%%%%%%%%%%%
%Typesetting conventions
%Lie groups and stuff
\def\gl#1#2{\ifmmode  {GL}(#1; {\bf #2}) \else $ {GL}(#1; {\bf #2})$\fi}
\def\sl#1#2{\ifmmode  {SL}(#1; {\bf #2}) \else $ {SL}(#1; {\bf #2})$\fi}
\def\so#1{\ifmmode  {SO}({#1}) \else $ {SO}(#1)$\fi}
	\newcommand{\sox}[2][R]{ $ {SO}(#2;{\bf #1})$}
	\newcommand{\soxx}[3][R]{ $ {SO}(#2,#3;{\bf #1})$}
\def\sp#1#2{\ifmmode  {Sp}(#1; {\bf #2}) \else $ {Sp}(#1; {\bf #2})$\fi}
\def\usp#1{\ifmmode  {USp}(#1) \else $ {USp}(#1)$\fi}
\def\spin#1{\ifmmode  {Spin}(#1) \else $ {Spin}(#1)$\fi} 
\def\su#1{\ifmmode  {SU}({#1}) \else $ {SU}(#1)$\fi}
%multiplet and dimension
\def\mult #1#2{\mathrm{#1}[\![ {#2}]\!]}
%normal ordering
\def\normal#1{\!:\! #1 \!:\! }
%doublebar
\def\double #1{#1{\hbox{\kern-2pt $#1$}}}
%underline
\def\un#1{\underline #1}
%Bras and kets
\def\ket#1{\left| #1\right>}
\def\bra#1{\left< #1\right|}
\def\braket#1#2{\left< #1\right|\left.#2\right>}
%For twistors
\def\kra#1{\left[ #1\right|}
\def\bet#1{\left| #1\right]}
%vev
\def\vev#1{{\langle #1\rangle}}
%small index
\def\nath#1{ {\mbox{\tiny{$(#1)$}}} }
%\def\nath#1{(#1)}

\def\unbr#1#2{\scriptsize{\mbox{$\underbrace{#1}_{#2}$}}}

\def\stack#1#2{\begin{array}{c}#1\\#2\end{array}}


%%%%% Local stuff %%%%
\def\S{{\cal S}}
\def\T{{\cal T}}
\def\U{{\cal U}}
\def\V{{\cal V}}
\def\W{{\cal W}}

%\usepackage{/Users/wdlinch3/Dropbox/Rashoumon/LaTeX/macros/WDL3}


\setbeamertemplate{bibliography entry title}{}
\setbeamertemplate{bibliography entry location}{}
\setbeamertemplate{bibliography entry note}{}

%%%%%%%%%%%%%%%%%%%%%%%%%%%%%%%%%%%%%%%%%%%%%%%%%%%
%%%%%%%%%%%%%%%%%%%%%%%%%%%%%%%%%%%%%%%%%%%%%%%%%%%


\usepackage{beamerthemesplit}

\title{Quantizable Models for M-theory}
%\author{William D. Linch III \\
%~\\
%{\em George P. and Cynthia W. Mitchell Institute for\\
%Fundamental Physics and Astronomy, \\
%Texas A\&{}M University.}
%}
\author{William D. Linch III} 
\date{
{\em George P. and Cynthia W. Mitchell Institute for\\
Fundamental Physics and Astronomy, \\
Texas A\&{}M University.}\\
~\\
7 September, 2015.\\
(Based on work with Warren Siegel \cite{Linch:2015lwa,Linch:2015fya,Linch:2015qva,Linch:2015fca}.)

}

\begin{document}

\frame{\titlepage}

\section[Outline]{}
%\frame{\tableofcontents}

%%%%%%%%%%%%%%%%%%%%%%%%%%%%%%%%%%%%%%%%%%%%%%%%%%%
\section{Review: {\bf S}trings, {\bf T}-duality, {\bf M}-theory, and {\bf F}-theory}
%%%%%%%%%%%%%%%%%%%%%%%%%
\frame{
\frametitle{S-theory}
\begin{itemize}
\item Pre-1984: There is a unique 11D supergravity and there are five 10D supergravities of types I, II{\sc a/b}, H{\sc e/o}.
%
\mypause
%
\item Green \&{} Schwarz (1984) \cite{Green:1984sg}:
10D supergravities ``$S(X)$'' on spacetime $X$ originate from quantization of fundamental strings ``$\mathbf S(X)$'' with target $X$. \\
%
\mypause
%
\item Moduli space of Type-II supergravity $S(X)$ on $X \to X\times T^n$:
\begin{align*}
{O(n,n) \over O(n)\times O(n)} 
\end{align*}
due to the symmetry between momentum modes (one copy of $n$) and winding modes (the other copy of $n$).
\item Narain (1986) \cite{Narain:1985jj}: $O(n,n)$ symmetry of $S(X)$ broken to $O(n,n;\mathbf Z)$ ``T-duality'' of the superstring $\mathbf S(X)$.
\end{itemize}
}

%%%%%%%%%%%%%%%%%%%%%%%%%
\frame{
\frametitle{T-theory}
\begin{itemize}
\item Duff (1989)$^{\scalebox{0.15}{\includegraphics{TAMUlogo.jpg}}}$
\cite{Duff:1989tf}: String theories have a generalization that manifests T-duality at the cost of doubling the spacetime dimension $x \to (x, \tilde x)$ for winding 0-modes.\\
%
\mypause
%
\item Siegel (1993) \cite{Siegel:1993xq,Siegel:1993th,Siegel:1993bj}: Doubled spacetime is reduced by a section condition on the momenta
\begin{align*}
(\partial_M) := \left(\begin{array}{c} \partial_\mu \\ \tilde \partial^\mu \end{array}\right)
~~~
\eta := \left(\begin{array}{cc} 0 & \delta_\mu^\nu \\ \delta^\mu_\nu & 0\end{array}\right)
~~~
\partial^M \partial_M = 0 
~~~\Rightarrow~~~\mathrm{e.g.} ~~~
\tilde \partial^m = 0.
\end{align*}
%
\mypause
%
Extension to a string theory ${\bf T}(X)$ with manifest T-duality symmetry.\\
%
\mypause
%
New Lie derivative (D-bracket) and commutator (C-bracket) described by current algebra of $O(D,D)$-covariant ``generalized'' momenta $\Pi_M = P_M + \eta_{MN}X^{\prime N}$ with quantum bracket 
\begin{align*}
[\Pi_M(\sigma_1) , \Pi_N(\sigma_2) ] = 2i \eta_{MN} \delta'(\sigma_1-\sigma_2).
\end{align*}
%
\mypause
%
Section condition arises dynamically as a new $O(D,D)$-invariant Virasoro constraint
\begin{align*}
L = \tfrac12 \eta^{MN} \Pi_M \Pi_N
~~~\Rightarrow~~~
\partial^M \partial_M = 0 .
\end{align*}
\end{itemize}
}

%%%%%%%%%%%%%%%%%%%%%%%%%
\frame{
\frametitle{M-theory}
\begin{itemize}
\item Cremmer, Julia, \& Scherk (1978) \cite{Cremmer:1978km}: Construction of eleven-dimensional supergravity. 
What about {\em its} fundamental origins?
%
\mypause
%
\item Bergshoeff, Sezgin, \& Townsend (1987) \cite{Bergshoeff:1987cm,Bergshoeff:1987qx}: No strings in 11D but super-membranes.
Strings do not explain existence of 11D supergravity but what are quantum membranes?
%
\mypause
%
\item Witten (1995) \cite{Witten:1995ex}: S-duality of the type-II{\sc a} $\bf S$tring theory has 11-dimensional supergravity as its massless truncation.
%
\begin{displaymath}
\xymatrix@=10pt{ 
&\ar@{-}@/^1pc/[dl]
	{\begin{array}{c}11\textrm{\sc d}\\ \textrm{{\sc sugra}}\end{array}} 
	\ar@{-}@/_1pc/[dr]& \\
\mathbf {II}\textrm{\sc a}\ar@{-}@/^1pc/[dd]&&\mathbf {H}\textrm{\sc e}\ar@{-}@/_1pc/[dd]\\
&\mathbf M&\\
\mathbf {II}\textrm{\sc b}&&\mathbf {H}\textrm{\sc o}\\
&\ar@{-}@/_1pc/[ul]{\mathbf I}\ar@{-}@/^1pc/[ur]&
}
\end{displaymath}
\end{itemize}
}

%%%%%%%%%%%%%%%%%%%%%%%%%
\frame{
\frametitle{11D Supergravity and U-duality}
Insights into the fundamental structure of M-theory?
%
\mypause
%
\begin{itemize}
\item Cremmer \&{} Julia (1979) \cite{Cremmer:1979up} 
and 
Hull \&{} Townsend (1994) \cite{Hull:1994ys}:
Analogously to the moduli space of compactifications of the type-II string theories on $n$-tori, hidden {\em exceptional} ``U-duality'' symmetry $E_{n(n)}$ of (gauged) 11D supergravity manifested by compactification on $T^n$.
\begin{table}[h]
\begin{center}
\begin{tabular}{|c|cccc|ccc}
\hline
$n$ &  $E_{n(n)}$  & $H_n$ & \# vectors & \# scalars \\
\hline
3 & \sl2R $\times$ \sl3R & \so2 $\times$ \so3 & 6 & 7 \\
4 & \sl5R & \so5  & 10 & 14\\
5 & \spin{5,5} & \so5 $\times$ \so5 & 16 & 25 \\
6 & $E_{6(6)}$ & $Sp(4)$ & 27 & 42 \\
7 & $E_{7(7)}$  & $SU(8)$ & 28 & 70\\
8 & $E_{8(8)}$ &  $SO(16)$ & 0 & 128 \\
\hline
\end{tabular} 
\end{center}
\end{table}
%
\mypause
%
\item Duff (1985) \cite{Duff:1985bv}: 
% at celebration of E. S. Fradkin
%(see also \cite{Duff:1986ne,Duff:1989tf}): 
``Either the $E_8\times SO(16)$ structure appears only in $d=3$ as an \underline{internal} symmetry, or else it is really present at every $d\leq 11$ as a \underline{spacetime} symmetry which transforms states of different spin into each other, in particular transforming [the frame field] into [the 3-form] in $d=11$. We emphasize that we are advocating the latter.'' 
\end{itemize}
}

%%%%%%%%%%%%%%%%%%%%%%%%%
\frame{
\frametitle{F-theory?}
Vafa (1996) \cite{Vafa:1996xn}: $\sl2Z$ S-duality of the type-II{\sc b} $\bf S$tring may be realized geometrically in ``12D'' theory $F$ 
\begin{align*}
\xymatrix{
	&F(\widehat X) \ar[dd] &\cr
	M( X\times S^1 )\ar@{<.}[ur] \ar[dr]&& \cr % T(X)\ar@{<.}[ul]_{\mathrm{reduction}}\ar[dl] 
	&S(X) &
	}
\end{align*}
%
\mypause
%
Defining properties of $F$:
\begin{enumerate}
\item Geometrically realized (S(trong/weak) part of) U-duality
\item Not really 12D: Type-II{\sc b} $\bf S$tring with axio-dilaton profile determined by elliptically-fibered Calabi-Yau 4-fold $T^2 \rightarrow \widehat X\rightarrow X$
\item Related by compactification to M-theory
\end{enumerate}
}

%%%%%%%%%%%%%%%%%%%%%%%%%
\frame{
\frametitle{Definition of $F(X)$} 
We propose to generalize the definition of F-theory by using both doubled and exceptional geometry \cite{Linch:2015lwa}: \\
%
\mypause
%
For any $X$, with supergravity $S(X)$, M-gravity $M(X)$, and generalized \& doubled gravity $T(X)$, define $F(X)$ to make the following diagram commute:\footnote{The rank $n=D+1$.}
\begin{align*}
\xymatrix{
	&{\stack{F(X)\vspace{1mm}}{E_{n(n)}/H_{n(n)} }} \ar[dl]
		\ar[dr] &\cr
{\stack{M(X)\vspace{1mm}}{ {GL({D+1},\mathbf R) \over O({D,1})} } }\ar[dr]_{\textrm{compactification~~~~~}}&& {\stack{T(X)\vspace{1mm}}{O({D,D})\over O(D)\times O(D) }}\ar[dl]^{\textrm{~~~sectioning}}  \cr
	&{\stack{S(X)\vspace{1mm}}{GL({D},\mathbf R) \over O({D-1,1})}} &
}
\end{align*}
}

%%%%%%%%%%%%%%%%%%%%%%%%%
\frame{
\frametitle{Extention to fundamental theories} 
Goal: For each rank $n=D+1$ find brane theory (generalizing the string), the quantization of which gives the theory $\mathbf F(X)$ that completes the diagram
\begin{align*}
\xymatrix{
	&\mathbf F(X) \ar[dl]_{\textrm{section~~~}} 
		\ar[dr]^{\textrm{~~~reduction}} &\cr
\mathbf M(X) \ar[dr]&& \mathbf T(X)\ar[dl] \cr
	&\mathbf S(X)&
}
\end{align*}
%
\mypause
%
It is {\em a priori} unclear whether this problem (existence of ${\bf F}$) has a solution. \\
In fact, we have the following 
}

%%%%%%%%%%%%%%%%%%%%%%%%%
\frame{
\frametitle{}
\begin{result*}[\cite{Linch:2015fya,Linch:2015qva,Linch:2015fca}]
Let $D= \mathrm{dim}(X)\leq 6$. %and let $n:=D+1$. 
Then ${\bf F}(X)$ exists in the form of a current algebra with exceptional global symmetry $G= E_{n(n)}$ of rank $n=D+1$. 
This current algebra arises from the quantization of a free theory with Hamiltonian symmetry $H_{n(n)}$ the (split form of the) maximal compact subgroup of $E_{n(n)}$.
When a Lagrangian formulation is known ($n\leq 5$), the symmetry is enhanced to $L_{n(n)}$ and the theory is that of a free self-dual form on a brane of split signature.
\end{result*}
{\renewcommand{\arraystretch}{1.3} %adds some padding
\begin{table}[h]
\begin{center}
\begin{tabular}{|c|ccc|}
\hline
$n$ 	& Lagrangian $L_{n(n)}$ & Hamiltonian $H_{n(n)}$ & Current Algebra $E_{n(n)}$ \\
\hline 
2	& \sl2R 	& \so{1,1} 	& \gl2R \\
3 	& $\sl2R^2\times\so2$ & $\sl2R\times\so2$ & $\sl3R\times\sl2R$\\
4 	& $\sl4R \simeq \spin{3,3}$ & $\sp4R\simeq\spin{3,2}$ & \sl5R \\
5 	& $\sl4C \simeq {Spin(6;{\bf C})}$	& \sp4C & \spin{5,5} \\
6 	& ?` $SU^\ast(8)$ ? & \usp{4,4} & $E_{6(6)}$ \\
7 & ?` ${U}^\ast(8)$ ? & \sl4H & $E_{7(7)}$\\
\hline
\end{tabular}
\end{center}
\end{table}
} %end arraystretch
}

%%%%%%%%%%%%%%%%%%%%%%%%%%%%%%%%%%%%%%%%%%%%%%%%%%
\section{Case Study: $\mathbf F(X)$ for the 3D String}
%%%%%%%%%%%%%%%%%%%%%%%%%
\frame{
\frametitle{Case Study: 3D String}
Properties of the F-gravity diamond for case with $X$ modeled on $\mathbf R^{2,1}$:
\begin{itemize}
%\item Focus on case with $X$ modeled on $\mathbf R^{2,1}$. %$D=\mathrm{dim}(X)=3$:\\
%
\mypause
%
\item Supergravity $S(X)$ is 3D, $N=2$ supergravity; Lorentz group is \so{2,1}.
%
\mypause
%
\item Manifestly T-dual generalized \& doubled SG has % duality group 
$\so{2,2} \sim \so{2,1}^2$.
%
\mypause
%
\item M-gravity $M(X)$ is 4D, $N=1$ ``old-minimal supergravity'' (with real prepotential for conformal compensator) \cite{Polacek:2013nla, Linch:2015lwa}.
%\footnote{Modified so that the chiral compensator has a real prepotential.} 
%
\mypause
%
\item F-gravity $F(X)$ has exceptional U-duality symmetry of rank 4: $G=E_{4(4)}\simeq \sl5R$.
%
\mypause
%
\item Coordinates $X^{mn}$ in the $\bf 10$ (signature 4+6).
%
\mypause
%
But solution of sectioning
\begin{align*}
\varepsilon^{mnpqr}\partial_{mn} \partial_{pq} = 0
~~~\Rightarrow ~~~
(\partial_{mn}) = (\partial_{-1\mu}, \partial_{\mu \nu}) \to (\partial_\mu, 0)
\end{align*}
reduces this to $(3+1)$-dimensional $M(X)$ breaking $\sl5R \to \gl4R$ spontaneously \cite{Berman:2012vc}.
\end{itemize}
}

%%%%%%%%%%%%%%%%%%%%%%%%%
\frame{
\frametitle{Virasoro currents} 
Strategy: Attempt to interpret strong constraint as a Virasoro-type constraint for the Hamiltonian analysis of some worldvolume theory.
\begin{itemize}
%
\mypause
%
\item In $\mathbf T(X)$ we saw the left-moving worldsheet current 
\begin{align*}
\Pi_M(\sigma) = P_M + \eta_{MN}X'^N
~~~\mathrm{with}~~~
[P_M(\sigma_1), X^N(\sigma_2)] = -i \delta_M^N \delta(\sigma_1-\sigma_2) 
\end{align*}
%\begin{displaymath}
%\xymatrix{
%	\Pi_\mu(\sigma)=&P_\mu(\sigma)+X'_\mu(\sigma) &;~~~ [P_\mu(\sigma_1), X^\nu(\sigma_2)] = -i \delta_\mu^\nu \delta(\sigma_1-\sigma_2) .\\
%\textrm{momentum} \ar[ur]  && \textrm{winding} \ar[ul]  &
%}
%\end{displaymath}
and Virasoro constraint $L =\tfrac12 \eta^{MN}\Pi_M\Pi_N$ giving sectioning $\partial^M\partial_M=0$.
%
\mypause
%
\item So we postulate for $\mathbf F(X)$ at each fixed $\tau$ a wordvolume with 5 coordinates $\sigma^m$ (in addition to the evolution coordinate $\tau$) and $G = \sl5R$-covariant current
\begin{align*}
\Pi_{mn} := P_{mn} + \tfrac12 \varepsilon_{mnpqr} \partial^r X^{pq}.
\end{align*}
%
\mypause
%
\item Virasoro-type generator for strong constraint on string fields:
\begin{align*}
\S^r := \tfrac1{16} \varepsilon^{mnpqr} \Pi_{mn}\Pi_{pq} .
\end{align*}
%
\mypause
%
\item Implies $\varepsilon^{mnpqr}(\partial_{mn}\, f)(\partial_{pq}\, g) = 0$ for all functions $f$ and $g$ on $F(X)$.
\end{itemize}
}

%%%%%%%%%%%%%%%%%%%%%%%%%
\frame{
\frametitle{Problem ... and solution}
But now algebra does not close:
\begin{align*}
[\S^m , \S^n] =   \S^{(m} \, \partial^{n)}\delta + \partial^{[m}\S^{n]} \,\delta + \varepsilon^{mnpqr} \Pi_{pq} \U_r \, \delta
~~~\mathrm{with}~~~
\U_r := \partial^s \Pi_{rs}.
\end{align*}
%
\mypause
%
$\Rightarrow$ Must interpret $\U$ as new constraint. It is central and reduces to 
\begin{align*}
\U_m = \partial^n P_{mn} 
%~~~...
\end{align*} 
%
\mypause
%
... but this is a Gau\ss{} law constraint! 
%
\mypause
%
Analogously to the first-quantization of Maxwell theory \cite{dirac2001lectures}: Quantize a 6D gauge 2-form $X_{MN}$!
{\renewcommand{\arraystretch}{1.3} %adds some padding
\begin{align*}
\hspace{-5mm}
\begin{array}{c|ccccccc}
	&\stack{\textrm{gauge}}{\textrm{field}} 
	&\stack{ \textrm{Lagrange} }{ \textrm{multiplier} } 	
	&\stack{\textrm{dynamical}}{\textrm{components}}
	&\stack{ \textrm{conjugate} }{ \textrm{momentum} }
	&\stack{\textrm{Gau\ss{}'s}}{\textrm{law}}
\\ \hline
\textrm{Maxwell} & A_\mu & A^\tau & A^i & E^i & \partial^i E_i
\\
\mathbf F(X) &X_{MN} & X^{\tau m} & X^{mn} & P_{mn} & \partial^nP_{mn} 
\end{array}
\end{align*}
}
}

%%%%%%%%%%%%%%%%%%%%%%%%%
\frame{
\frametitle{Chiral forms}
Siegel (1983) \cite{Siegel:1983es}: 
By analogy with string action for chiral scalar fields, 
action for gauge 2-form with self-dual fieldstrength $F_{MNP}$:
\begin{align*}
S_{iegel} &= \int 
	\left(
	\tfrac 1{12} F^{MNP} F_{MNP} + \tfrac14 \lambda^{MN} F_M^{(-)PQ}F^{(-)}_{NPQ}
	\right)\,d^6\sigma  .	
\end{align*}
%
\mypause
%
The equation of motion for $\lambda$
\begin{align*}
F_M^{(-)PQ}F^{(-)}_{NPQ} = 0
~~~\Rightarrow~~~
F^{(-)}_{MNP} = 0	
\end{align*}
implies the fieldstrength of $X$ is self-dual.
%
\mypause
%
The gauge transformations
\begin{align*}
\delta X_{MN} &= \alpha^P F^{(-)}_{MNP} + \tfrac1{12} \alpha^P \lambda_{[M}{}^QF{(-)}_{NPQ]} 
	+\mathcal O (\lambda^2) 
\cr
\delta \lambda^{MN} &= \partial^{(M} \alpha^{N)}
	+ \alpha^P\partial_P \lambda^{MN}
	+ \tfrac12 \lambda_P{}^{(M} \left(
		\partial^{N)} \alpha^P - \partial^P \alpha^{N)} 
		\right)
-\textrm{trace}		
+\mathcal O (\lambda^2) 
\nonumber	
\end{align*}
imply $\lambda$ is pure gauge.
}

%%%%%%%%%%%%%%%%%%%%%%%%%
\frame{
\frametitle{Canonical analysis}
In  $\lambda \to 0$ gauge the Hamiltonian action becomes
\begin{align*}
S &= 
	-\int  \tfrac12 \dt X{}^{mn} P_{mn} \,d^5\sigma d\tau
	+\int  H \, d\tau
\cr
H&= \int 
	\left(
	\tfrac14 P^{mn} P_{mn} 
	+ \tfrac1{12} F^{mnp}F_{mnp}
	+ X^{\tau m} \U_m 
	\right) \, d^5 \sigma .	
\end{align*}
%
\mypause
%
Current from selfdual fieldstrength:
\begin{align*}
\Pi_{mn} := F^{(+)}_{\tau mn} = P_{mn} +\tfrac12 \varepsilon_{mnpqr} \partial^r X^{pq} .
\end{align*}
%
\mypause
%
Free quantum brackets $[P_{mn}(\sigma_1) , X^{pq}(\sigma_2) ] = -i \delta_{[m}^{[p}\delta_{n]}^{q]} \delta(\sigma_1 - \sigma_2)$ imply
\begin{align*}
\boxed{[\Pi_{mn}(\sigma_1) , \Pi_{pq}(\sigma_2) ] = 2i \epsilon_{mnpqr} \partial^r \delta(\sigma_1 - \sigma_2) }
\end{align*}
%
\mypause
%
Stress-energy tensor
\begin{align*}
\T_{mn} = \eta_{mn} \T - \tfrac12 \eta^{pq}\Pi_{mp}\Pi_{nq}
~~~
\boxed{\S^r = \tfrac1{16} \epsilon^{mnpqr} \Pi_{mn} \Pi_{pq}}
~~~
\T = \tfrac18 \eta^{mp}\eta^{nq} \Pi_{mn}\Pi_{pq}
\end{align*}
%
\mypause
%
$\Rightarrow$ %~~~~Existence of c
Closed $E_{4(4)} = \sl5R$-covariant sub-algebra generated by  $\S$ and $\U$!
}

%%%%%%%%%%%%%%%%%%%%%%%%%
\frame{
\frametitle{Courant bracket} 
Spacetime gauge transformations generated by vector field currents
\begin{align*}
\Lambda_i := \tfrac i2 \int  \lambda_i^{mn}(X(\sigma)) \Pi_{mn}(\sigma) \,d^5\sigma .
\end{align*}
%
\mypause
%
Algebra
\begin{align*}
[\Lambda_1, \Lambda_2 ] = \Lambda_{12}
~~~\mathrm{with}~~~
\lambda_{12}^{mn} 
	&= \tfrac14 \lambda_{[1}^{pq}\partial_{pq} \lambda_{2]}^{mn}
	-\tfrac12 \lambda_{[1}^{p[m}\partial_{pq} \lambda_{2]}^{n]q}
	-\tfrac14 \lambda_{[1}^{mn}\partial_{pq} \lambda_{2]}^{pq} 
\end{align*}
exactly reproduces the ``exceptional Courant bracket'' of diffeomorphisms of the $E_{4(4)}$ F-gravity computed by Berman, Godazgar,${}^2$ and Perry \cite{Berman:2011cg}.\\
~\\
%
\mypause
%
Conclusion:\begin{enumerate}
\item Quantization of the chiral 2-form on a 5-brane with split signature gives rise to a current algebra with $E_{4(4)}$ symmetry. 
%
\mypause
%
\item The massless truncation of its C-bracket reproduces the ``exceptional Courant bracket'' of diffeomorphisms of the M-gravity appropriate to this case.
%
\mypause
%
\item Solution of the constraints reduces the current algebra to that of the three-dimensional bosonic string. 
\end{enumerate}
}

%%%%%%%%%%%%%%%%%%%%%%%%%%%%%%%%%%%%%%%%%%%%%%%%%%
\section{Generalizations}
%%%%%%%%%%%%%%%%%%%%%%%%%
\frame{
\frametitle{U-duality} 
The U-duality $E_{4(4)}$ can be enlarged to $G=E_{n(n)}$. \\
%
\mypause
%
Consider the Dynkin diagram
\begin{align*}
\put(-60,-3){$R_2$}
\multiput(-40,0)(20,0){5}{\circle{6}}
\multiput(-37,0)(20,0){2}{\line(1,0){14}}
\put(6,-0.5){...}
\put(23,0){\line(1,0){14}}
%\multiput(3,0)(20,0){2}{\line(1,0){14}}
\put(0,20){\circle{6}}
\put(-5,27){$R_3$}
\put(0,3){\line(0,1){14}}
\put(47,-3){$R_1$}
\end{align*}
%
\mypause
%
(This corresponds to $R_1 = \mathbf{10}^\prime$, $R_2 = \bf 5$, and $R_3 = \bf 5'$ for $E_{4(4)}= \sl5R$.)\\
%
\mypause
%
Let ${\bf m}$ run over $R_1)$ 
and 
$m$ run over $R_2$ and introduce the $(R_1 \otimes R_1)_{sym} \supset R_2$ Clebsch-Gordan coefficients $\eta_{{\bf mn}p}$ and $\eta^{{\bf mn}p}$.\\
%
\mypause
%
Generalization of the current algebra:
\begin{align*}
[\Pi_{\bf m}(\sigma_1) , \Pi_{\bf n}(\sigma_2) ] = 2i \eta_{{\bf mn} p} \partial^p \delta(\sigma_1, \sigma_2)
\end{align*}
%
\mypause
%
Virasoro constraint:
\begin{align*}
\S^p= \tfrac14 \eta^{{\bf mn} p} \Pi_{\bf m}\Pi_{\bf n} .
\end{align*}
}

%%%%%%%%%%%%%%%%%%%%%%%%%
\frame{
\frametitle{}
C-bracket of two vector fields:
\begin{align*}
[V_1^{\bf m} \Pi_{\bf m}(\sigma_1) , V_2^{\bf n} \Pi_{\bf n}(\sigma_1) ]
	&=
	2i \eta_{{\bf mn}p}\partial^p \delta(\sigma_1-\sigma_2) V_1^{\bf m}V_2^{\bf n}
\cr
	&-i \delta(\sigma_1-\sigma_2) 
		\left[
			V_{[1}^{\bf m}\partial_{\bf m}V_{2]}^{\bf n}
			- \tfrac12 \eta^{{\bf mn}r}\eta_{{\bf pq}r} V_1^{\bf p}\partial_{\bf m}V_2^{\bf q}
		\right] \Pi_{\bf n} .
\end{align*}
%
\mypause
%
Truncation to 0-modes should reproduce the $E_{n(n)}$ ``Courant brackets'' 
computed by
Berman, Cederwall, Kleinschmidt, and Thompson \cite{Berman:2012vc}.\\
%
\mypause
%
Therefore, the projector $(R_1 \otimes R_1)_{sym}\to R_2$ is identified with the Clebch-Gordan coefficients:
\begin{align*}
Y^{\bf mn}{}_{\bf pq} \leftrightarrow \eta^{{\bf mn}r}\eta_{{\bf pq}r} 
\end{align*}
up to constraints.\\ 
%
\mypause
%
This $Y$ projector factorizes for all ranks $n\leq 6$ and factorizes up to a new constraint for $E_{7(7)}$ (corresponding to the 6D string).\\
%
\mypause
%
Ultimate goal is 10-dimensional Lorentz invariance for $X$ corresponding to $E_{11}$ U-duality conjectured by West (2001) \cite{West:2001as}.
}

%%%%%%%%%%%%%%%%%%%%%%%%%%
%\frame{
%\frametitle{Critical superstrings} 
%Critical superstrings can be constructed by adding in $10-D$ scalars $Y$ (and their duals $\Tilde Y$) in the ``ordinary string way''. The Hamiltonian actions will have symmetry $H \times H'$ with $H'=\so{10-D}$. \\
%Fermions are added for supersymmetry which still looks like 
%\begin{align*}
%\{ D_{\bm \alpha}(\sigma_1) , D_{\bm \beta}(\sigma_2) \} = 2 (\gamma^{\hat {\bf m}})_{\bm \alpha \bm \beta} \Pi_{\hat {\bf m}} \delta(\sigma_1-\sigma_2)
%\end{align*}
%but now the Clifford algebra is deformed to a Clifford-ish algebra
%\begin{align*}
%\gamma^{\hat {\bf m}}\gamma^{\hat {\bf n}} + \gamma^{\hat {\bf n}}\gamma^{\hat {\bf m}}
%= 2 \eta^{\hat {\bf m} \hat{\bf n}} {\bf 1} + 2 \eta^{\hat {\bf m} \hat{\bf n} p} \gamma_p.
%\end{align*}
%Jacobi identities require the Fierz-ish identity
%\begin{align*}
%\eta_{\hat {\bf m} \hat{\bf n} p } (\gamma^{\hat {\bf m}})_{(\bm \alpha \bm \beta} (\gamma^{\hat {\bf n}})_{\bm \gamma) \bm \delta} = 0
%\end{align*}
%restricting $D+D' = 3, 4, 6, 10$.\footnote{Also provides twistor solution $p_{\bf m} = \lambda \gamma_{\bf m} \lambda$ to the section condition $\S=0$ modulo ``pure spinors'' $\lambda \gamma_{\bf m} \lambda = 0$.
%} %end footnote
%
%}

%%%%%%%%%%%%%%%%%%%%%%%%%%%%%%%%%%%%%%%%%%%%%%%%
%\section{Relation to Dirac-Born-Infeld Branes}
%\frame{
%\frametitle{Dirac-Born-Infeld current algebra}
%\begin{enumerate}
%\item Pasti, Sorokin, \& Tonin (1997) \cite{Pasti:1997gx}: Standard immersion $\phi: \Sigma \to X$ of probe M2 and M5 into 11D target space.
%\item Hatsuda \& Kamimura (2012): ``Exceptional'' symmetry $E_{4(4)}$ \cite{Hatsuda:2012vm} and $E_{5(5)}$ \cite{Hatsuda:2013dya} from M2 and M5 DBI actions in with target with dimension $\mathrm{dim}(\rho_2)$.
%\end{enumerate}
%}

%%%%%%%%%%%%%%%%%%%%%%%%%%%%%%%%%%%%%%%%%%%%%%%%
%\section{Retrospective}
%\frame{
%\frametitle{Retrospective}
%Why no quantizable membrane/higher-brane theories?\\
%Either gapless \cite{Polyakov:1986cs,Curtright:1982gt,Lovelace:1971fa} or no massless states at all \cite{deWit:1988ig,Mezincescu:1987kj}. \\
%Novelties and possible advantages of this version of $\bf F$:
%\begin{enumerate}
%\item Extended geometry of target $F(X)$ allows for more manifest U-duality.
%\item Manifest U-duality grows as we increase the Lorentz invariance of $X$ (Duff's conjecture).
%%\item Reduction of na\"ive target space dimension of $F(X)$ to $M(X)$ is dynamical, not a by-hand truncation
%\item Reduction to the {\em stable} (critical super)string is {\em dynamical} instead of a by-hand truncation (same degrees of freedom as usual strings, just more 0-modes for ``winding'').
%\item Extended geometry allows for conformal worldvolume theory: Instead of worldvolume scalar ``immersion map'' $\phi : \Sigma \to X$ with complicated dynamics, worldvolume field is a free self-dual $p$-form.
%\end{enumerate}
%}

%%%%%%%%%%%%%%%%%%%%%%%%%%%%%%%%%%%%%%%%%%%%%%%
\section{Future work}
%%%%%%%%%%%%%%%%%%%%%%%%%
\frame{
\frametitle{Open questions and current \& future work}
Loads of work to do... 
\begin{enumerate}
%
\mypause
%
\item Relation to M2 and M5 in 11D target space of Bergshoeff, Sezgin, \& Townsend (1987) \cite{Bergshoeff:1987cm,Bergshoeff:1987qx} and Howe \& Sezgin (1996) \cite{Howe:1996mx,Howe:1996yn}?
%
\mypause
%
\item More U-duality? (Can we use Palmkvist \cite{Palmkvist:2015dea} and related work to get $E_{>7}$?)
%
\mypause
%
\item Narian-type description in terms of exceptional Borcherds algebras? 
\begin{align*}
\put(-55,-3){$\mathcal S$}
\multiput(-40,0)(20,0){5}{\circle{6}}
\multiput(-37,0)(20,0){2}{\line(1,0){14}}
\put(6,-0.5){...}
\put(23,0){\line(1,0){14}}
\put(0,20){\circle{6}}
\put(-5,27){$\mathcal U$}
\put(0,3){\line(0,1){14}}
\put(47,-3){$P$}
\end{align*}
%
\mypause
%
\item Relation to (gravitational) tensor hierarchies?
%
\mypause
%
\item Worldvolume gravity? 
%
\mypause
%
\item Quantization?
\end{enumerate}
%
\mypause
%
\begin{center}
{\bf \huge THANK YOU!}
\end{center}
}

%%%%%%%%%%%%%%%%%%%%%%%%%%%%%%%%%%%%%%%%%%%%%%%
{\tiny
%\small
%\footnotesize
%\linespread{1.1}\selectfont
\bibliography{/Users/wdlinch3/Dropbox/Rashoumon/LaTeX/BibTex/BibTex}
\bibliographystyle{unsrt}
} % end resize


\end{document}